% !TeX root = ../初等数学笔记.tex
% 由主文档通过 \input 载入；请编译主文档。
\setcounter{chapter}{9} % 保留原书章号。
\chapter{初等不等式与最值}
\label{chap:em-inequality-methods}
不等式把精确求值转化为范围控制，最值问题则进一步要求找到最好的界。有效证明必须同时交代变量范围、每步估计依据与取等条件，否则即使得到了正确的界，也可能没有解决最值问题。

\section{作差、作商与平方非负}
证明 $A\ge B$ 可以化为 $A-B\ge0$；若 $B>0$，也可以证明 $A/B\ge1$。最基本的非负结构是平方，例如
\[
 a^2+b^2-2ab=(a-b)^2\ge0,
\]
等号当且仅当 $a=b$。含参数的二次式常以配方处理，几个变量的式子则尝试写成若干平方的和。例如
\[
 a^2+b^2+c^2-ab-bc-ca
 =\tfrac12\bigl((a-b)^2+(b-c)^2+(c-a)^2\bigr)\ge0,
\]
等号当且仅当 $a=b=c$。给不同项分别估计之后，必须检查这些等号是否能够同时成立。

\section{常见平均数与均值不等式}
对正数 $x_1,\ldots,x_n$，定义调和平均（harmonic mean）、几何平均（geometric mean）、算术平均（arithmetic mean）与平方平均（quadratic mean）：
\[
 H=\frac n{\sum_i1/x_i},\quad G=\left(\prod_i x_i\right)^{1/n},
 \quad A=\frac1n\sum_i x_i,\quad Q=\sqrt{\frac1n\sum_i x_i^2}.
\]
\begin{theorem}[均值链]
$H\le G\le A\le Q$，任一相邻等号成立当且仅当所有 $x_i$ 相等。
\end{theorem}
\begin{proof}
二元关系 $\sqrt{uv}\le(u+v)/2$ 来自 $(\sqrt u-\sqrt v)^2\ge0$。反复两两分组可证明个数为 $2^k$ 时几何平均不超过算术平均。对一般 $n$，选择 $m=2^k\ge n$，在原数后补入 $m-n$ 个 $A$，新算术平均仍为 $A$，所以 $(\prod_i x_i)A^{m-n}\le A^m$，除以正数得到 $G\le A$。反复分组的等号要求全部数相等，给出取等条件。对倒数应用同一结论得 $1/G\le1/H$，即 $H\le G$。最后 $\sum_i(x_i-A)^2=\sum_i x_i^2-nA^2\ge0$，故 $A\le Q$，等号也要求每项等于 $A$。
\end{proof}
$G\le A$ 可推广到非负数，含零时直接判断即可；$H$ 仍需所有数为正。相同距离分别以速度 $u,v>0$ 行驶，平均速度为 $2/(1/u+1/v)$；相同时间行驶则为 $(u+v)/2$。选择平均数取决于总量如何累积。

\section{柯西--施瓦茨不等式与分式估计}
\begin{theorem}[柯西--施瓦茨不等式]
对实数 $a_i,b_i$（$1\le i\le n$），
\begin{equation}\label{eq:em-cauchy}
 \left(\sum_i a_ib_i\right)^2\le
 \left(\sum_i a_i^2\right)\left(\sum_i b_i^2\right).
\end{equation}
等号当且仅当对任意 $i,j$ 都有 $a_ib_j=a_jb_i$；若 $(b_i)$ 不全为零，等价于存在实数 $\lambda$ 使 $a_i=\lambda b_i$。
\end{theorem}
\begin{proof}
直接展开得到左、右两端之差的非负表达：
\[
 \left(\sum_i a_i^2\right)\left(\sum_i b_i^2\right)
 -\left(\sum_i a_ib_i\right)^2
 =\sum_{i<j}(a_ib_j-a_jb_i)^2.
\]
各平方同时为零即为等号条件。若某个 $b_j\ne0$，取 $\lambda=a_j/b_j$ 即得比例形式；若全部 $b_i=0$，两端同为零。
\end{proof}
对 $y_i>0$，代入 $a_i=x_i/\sqrt{y_i},b_i=\sqrt{y_i}$ 得
\[
 \sum_i\frac{x_i^2}{y_i}\ge\frac{(\sum_i x_i)^2}{\sum_i y_i},
\]
等号当且仅当各 $x_i/y_i$ 相同。该形式适合分母不同但总和可控的分式估计。

\section{排序、伯努利与凸性选读}
\emph{排序不等式}（rearrangement inequality）断言：若 $a_1\le\cdots\le a_n$、$b_1\le\cdots\le b_n$，则对任意排列 $\sigma$，
\[
 \sum_i a_ib_{n+1-i}\le\sum_i a_ib_{\sigma(i)}\le\sum_i a_ib_i.
\]
证明右侧只需消除逆序：交换 $i<j$ 上逆序的两个 $b$ 后，和的增加量为 $(a_j-a_i)(b_{\sigma(i)}-b_{\sigma(j)})\ge0$；有限次排序得到最大值，反向排列同理得到最小值。与最大值取等当且仅当不存在 $a_i<a_j$ 而所配 $b$ 值严格逆序的数对；两列均严格递增时，最大值只能由同序排列取得。

对整数 $n\ge1$、实数 $x\ge-1$，\emph{伯努利不等式}（Bernoulli's inequality）为 $(1+x)^n\ge1+nx$。归纳时乘以 $1+x\ge0$，再用
$(1+nx)(1+x)=1+(n+1)x+nx^2\ge1+(n+1)x$。$n=1$ 恒取等，$n>1$ 时仅 $x=0$ 取等。

区间 $I$ 上的函数 $f$ 称为\emph{凸函数}（convex function），若任意 $u,v\in I$、$0\le t\le1$ 都满足
$f((1-t)u+tv)\le(1-t)f(u)+tf(v)$。反复合并两个权重可归纳得到有限 Jensen 不等式
\[
 f\left(\sum_i\lambda_i x_i\right)\le\sum_i\lambda_i f(x_i),
 \quad \lambda_i\ge0,\quad\sum_i\lambda_i=1.
\]
归纳步骤先去掉零权重，再把前 $n-1$ 项归一化合为一个点，先后应用两点凸性与归纳假设。若两点不同且 $0<t<1$ 时严格不等，称严格凸；此时取等当且仅当正权重对应的 $x_i$ 全相等。$f(x)=x^2$ 的凸性直接由差值 $t(1-t)(u-v)^2$ 验证，无需微分法。

\section{约束最值与等号的可达性}
设矩形周长为 $L>0$，边长 $x,y>0$，则 $x+y=L/2$。均值不等式给出 $xy\le((x+y)/2)^2=L^2/16$；$x=y=L/4$ 满足约束且取等，所以它确为最大面积。

若约束改为 $x>1$，研究 $x+1/x$，虽有下界 $2$，但取等要求 $x=1$，不在可行域。取 $x=1+\varepsilon$ 并令正数 $\varepsilon$ 越来越小，可使函数值任意接近 $2$，因此 $2$ 是最大下界，却不是最小值。这解释了“得到一个界”和“证明最值”之间所差的可达性检验。

\section{例题、习题与学习检查}
\begin{example}
已知 $x,y>0$ 且 $x+y=1$，求 $1/x+1/y$ 的最小值。
\end{example}
\begin{solve}
由 $xy\le1/4$，有 $1/x+1/y=(x+y)/(xy)\ge4$；$x=y=1/2$ 时取等。也可直接使用分式柯西不等式得到 $1/x+1/y\ge(1+1)^2/(x+y)=4$，两种证明指向同一平衡条件。
\end{solve}
\begin{enumerate}
\item 若 $x,y\ge0,x+y=6$，求 $xy$ 最大值。答：$9$，在 $x=y=3$ 取得。
\item 证明 $(a+b+c)^2\le3(a^2+b^2+c^2)$ 并给出等号条件。答：应用式 \eqref{eq:em-cauchy}，等号为 $a=b=c$。
\item 用二点凸性的定义证明 $f(x)=|x|$ 为凸函数，并说明它不是严格凸函数。
\end{enumerate}

